% ------------------------------------------------------------------------------
% Fundamentação Teórica
% ------------------------------------------------------------------------------

\chapter{Fundamentação Teórica}
\label{chap_fundamentacao_teorica}


Esta seção apresenta uma revisão dos fundamentos físicos e matemáticos da queda livre, abordando desde o modelo clássico de partícula até modelos para corpos rígidos com diferentes regimes de resistência do ar. São discutidas as principais equações, hipóteses e limitações de cada abordagem, com base em literatura científica \cite{galilei1638,halliday2011,batchelor2000,tipler2009,butcher2016}.

\section{Queda Livre de uma Partícula (Sem Resistência do Ar)}

No modelo clássico, considera-se uma partícula ideal em queda sob ação exclusiva da gravidade, desconsiderando qualquer resistência do ar. As equações do movimento são \cite{halliday2011,tipler2009}:
\begin{equation}
    a(t) = -g
    \label{eq:acel_queda_livre}
\end{equation}
\begin{equation}
    v(t) = v_0 - g t
    \label{eq:vel_queda_livre}
\end{equation}
\begin{equation}
    y(t) = y_0 + v_0 t - \frac{1}{2} g t^2
    \label{eq:pos_queda_livre}
\end{equation}
onde $g$ é a aceleração da gravidade, $v_0$ a velocidade inicial e $y_0$ a posição inicial. Essas equações são válidas para corpos pequenos ou densos, onde o efeito do ar é desprezível \cite{halliday2011,tipler2009}.



\section{Queda Livre de Corpos Rígidos}

Ao analisar a queda livre de um corpo rígido, é importante considerar fatores como área frontal, distribuição de massa e possíveis rotações. No vácuo, o centro de massa do corpo rígido segue as mesmas equações da partícula, com aceleração constante igual à gravidade. No entanto, para corpos extensos, pode-se discutir limitações do modelo de partícula e a adequação do modelo para situações reais \cite{halliday2011,tipler2009}.

Quando a resistência do ar não pode ser desprezada, dois regimes principais são considerados:

\begin{itemize}
    \item \textbf{Resistência linear (Lei de Stokes):} Válida para baixas velocidades e partículas pequenas, o arrasto é proporcional à velocidade \cite{batchelor2000,butcher2016}:
    \begin{equation}
        F_d = -6 \pi \eta r v
    \end{equation}
    \begin{equation}
        m \frac{dv}{dt} = -mg - 6\pi\eta r v
    \end{equation}
    onde $\eta$ é a viscosidade dinâmica do ar e $r$ o raio da esfera \cite{batchelor2000,butcher2016}. O termo $6\pi\eta r$ é chamado de $c$ (ou $k$ em alguns textos):
    \begin{equation}
        c = 6\pi\eta r
    \end{equation}
    Assim, a equação do movimento pode ser reescrita como:
    \begin{equation}
        m \frac{dv}{dt} = -mg - c v
    \end{equation}
    Essa notação facilita a identificação do parâmetro de resistência linear nos experimentos e simulações.
    
    \item \textbf{Resistência quadrática:} Para velocidades mais altas ou corpos maiores, o arrasto é proporcional ao quadrado da velocidade \cite{batchelor2000,butcher2016}:
    \begin{equation}
        F_d = -\frac{1}{2} \rho C_d A v^2
    \end{equation}
    \begin{equation}
        m \frac{dv}{dt} = -mg - \frac{1}{2} \rho C_d A v^2
    \end{equation}
    onde $\rho$ é a densidade do ar, $C_d$ o coeficiente de arrasto e $A$ a área frontal do corpo \cite{batchelor2000}. O termo $\frac{1}{2} \rho C_d A$ é chamado de $k$ (ou $c$ em alguns textos):
    \begin{equation}
        k = \frac{1}{2} \rho C_d A
    \end{equation}
    Assim, a equação do movimento pode ser reescrita como:
    \begin{equation}
        m \frac{dv}{dt} = -mg - k v^2
    \end{equation}
    Essa notação facilita a identificação do parâmetro de resistência quadrática nos experimentos e simulações.
\end{itemize}

Esses modelos permitem analisar o comportamento do sistema em diferentes regimes, avaliar a influência da resistência do ar sobre a trajetória, a velocidade terminal e a acurácia dos métodos numéricos, além de comparar os efeitos de cada abordagem para situações práticas.

% ------------------------------------------------------------------------------
\section{Considerações sobre a Resistência do Ar}


A resistência do ar é um fator fundamental na modelagem do movimento de corpos em queda livre. Sua influência depende do tamanho, forma, densidade do corpo e do regime de velocidade \cite{halliday2011,batchelor2000}. Para objetos pequenos, densos e em trajetórias curtas, a resistência do ar pode ser desprezada, justificando o uso do modelo clássico de partícula. No entanto, para corpos maiores, menos densos ou em velocidades elevadas, a resistência do ar se torna significativa e deve ser incorporada ao modelo físico. Os dois principais regimes de resistência são:

\begin{itemize}
    \item \textbf{Resistência linear:} Predomina em baixas velocidades e para partículas pequenas, sendo proporcional à velocidade (Lei de Stokes) \cite{batchelor2000}.
    \item \textbf{Resistência quadrática:} Torna-se relevante em velocidades mais altas e para corpos maiores, sendo proporcional ao quadrado da velocidade (modelo aerodinâmico completo) \cite{batchelor2000,butcher2016}.
\end{itemize}

A escolha do modelo adequado é essencial para garantir a precisão das simulações e a correta interpretação dos resultados experimentais.

A fundamentação teórica apresentada evidencia que a modelagem da queda livre pode ser adaptada conforme o grau de realismo desejado e os objetivos do estudo. O caso clássico, sem resistência do ar, fornece uma base sólida para a compreensão dos conceitos fundamentais e permite validação direta dos métodos numéricos. A inclusão de um corpo rígido e, principalmente, da resistência do ar — seja no regime linear ou quadrático — torna o modelo mais próximo da realidade, permitindo analisar situações práticas e discutir as limitações das hipóteses simplificadoras.

Além disso, a análise dos diferentes regimes de resistência e a avaliação da acurácia dos métodos numéricos em função do passo de tempo são essenciais para garantir resultados confiáveis e interpretar corretamente os fenômenos físicos simulados. Assim, a fundamentação teórica não só embasa a implementação computacional do modelo de queda livre, como também orienta a escolha das abordagens mais adequadas para cada situação, em consonância com os objetivos propostos pelo roteiro do trabalho.


